\section{Introduction}
\label{intro}
Retrieval-Augmented Generation (RAG) has become a standard approach for leveraging Large Language Models (LLMs) in enterprise information tasks. By feeding retrieved documents into the generation process, RAG addresses the static nature of parametric memory in pretrained models and significantly reduces hallucinations, as discussed in the paper \textit{Retrieval-Augmented Generation for Knowledge-Intensive NLP Tasks} \parencite{lewis2020retrieval}. However, this reliance on external information also introduces a new attack surface. Specifically, if the retrieved context is corrupted, poisoned, or semantically altered, the model has to deal with a conflict between what it already knows and what it is being told, which can compromise the accuracy and reliability of its outputs.

Recent work has started to look into this problem, showing that RAG systems tend to break down when faced with misinformation. In theory, a robust RAG system should be able to detect when the retrieved context is contradictory or clearly wrong and refuse to answer (abstain). In practice, however, LLMs often exhibit what has been called the ``Factual Accuracy Paradox'', as presented in the paper \textit{The Power of Noise: Redefining Retrieval for RAG Systems} \parencite{cuconasu2024power}, where irrelevant or noisy documents seriously degrade the model's ability to produce truthful answers.

Although prior studies have established that this vulnerability exists, several important questions remain open. In particular, there has been limited analysis of how the severity of this vulnerability changes with model size (scaling laws) and whether it can be mitigated through prompt design (Prompt Engineering). On top of that, evaluating these vulnerabilities at scale is itself a challenge. Manual human evaluation is time-consuming, expensive, and hard to scale, which is why LLM-as-a-Judge methods have been gaining traction. But the reliability of these methods is still debated, especially when the task requires distinguishing between a model that simply parrots back wrong information from retrieved documents and one that uses its own knowledge to correct and give a better answer.

This study uses Semantic Mutation Testing to evaluate RAG systems. The goal is to analyze how model size (Llama 3.1 8B and 17B parameters), the use of system prompts, and whether the data is public or private affect the system's ability to resist semantically poisoned documents. To support evaluation at scale, we built an automated evaluation pipeline based on a Mixture of Experts (MoE) architecture. In this pipeline, Llama 4 Scout 17B serves as the AI Judge, followed by a human evaluation step. This setup allows us to not only assess model performance but also examine how reliable the LLM-as-a-Judge approach really is in practice.
